
\documentclass[11pt,a4paper]{article}

% ============================================================
% VANTA: Preliminary Research Manuscript
% Version 0.1 — November 2026
% ============================================================

\usepackage[margin=0.95in]{geometry}
\usepackage{amsmath,amssymb,amsthm}
\usepackage{booktabs}
\usepackage{graphicx}
\usepackage{hyperref}
\usepackage{microtype}
\usepackage{enumitem}

\hypersetup{
    colorlinks=true,
    linkcolor=blue,
    citecolor=blue,
    urlcolor=blue
}

\newtheorem{proposition}{Proposition}
\newtheorem{remark}{Remark}

\newcommand{\R}{\mathbb{R}}
\newcommand{\od}{\odot}

% ============================================================
% Title
% ============================================================

\title{
    \textbf{VANTA: Value-Aggregated Non-Softmax Token Attention}\\
    \large A Lightweight Feature-Interaction Architecture
    for Tabular Regression
}

\author{
    Parv Sharma\\
    \textit{Preliminary Research Manuscript}
}

\date{November 2026}

\begin{document}

\maketitle

% ============================================================
% Abstract
% ============================================================

\begin{abstract}

We investigate whether learned feature-wise attention
can provide an effective inductive bias for tabular
regression without requiring a conventional Transformer
architecture.

We introduce VANTA (Value-Aggregated Non-Softmax
Token Attention), a lightweight regression architecture
that constructs asymmetric feature-interaction scores
using learned affine query and key projections.

Unlike conventional softmax attention, VANTA retains
signed, unnormalized dot-product interactions and
aggregates them using scalar feature values directly.
The resulting representation is combined with the
original feature space through a residual connection
and mapped to a scalar prediction using a linear
readout.

We examine two parameterizations: VANTA-R, which
uses a learned scalar residual gate, and VANTA-G,
which uses feature-wise residual gates.

We derive an associative formulation requiring
$\mathcal{O}(Fd)$ arithmetic for feature aggregation,
where $F$ is the number of input features and $d$
is the latent interaction dimension.

On the Diamonds regression dataset, VANTA-R
achieves a test $R^2$ of 0.9740 using 1,450
trainable parameters, while VANTA-G achieves
0.9742 using 1,458 parameters.

Ablation experiments suggest that learned query-key
interactions contribute substantially to predictive
performance. However, stronger nonlinear baselines,
including a parameter-matched multilayer perceptron
and gradient-boosted trees, achieve higher accuracy
on the same benchmark.

These preliminary results motivate further
investigation into the generalization properties
of low-rank asymmetric feature-interaction
operators for tabular regression.

\end{abstract}

\noindent\textbf{Keywords:}
Tabular regression, feature interactions,
attention, low-rank factorization,
residual learning, categorical embeddings.

% ============================================================
\section{Introduction}
% ============================================================

Tabular regression seeks to estimate a continuous
target variable from a heterogeneous collection
of numerical and categorical features.

Linear regression provides a simple and interpretable
baseline, but cannot directly represent nonlinear
feature interactions without explicitly constructed
interaction terms.

Neural networks can approximate nonlinear functions,
although their increased flexibility may require
substantial data and careful regularization.

Attention-based architectures provide an alternative
mechanism for learning interactions between input
representations. However, conventional Transformer
architectures combine attention with additional
components, including deep feed-forward networks,
multiple attention layers, and normalization modules.

For relatively small tabular problems, these
components may introduce unnecessary complexity.

This work investigates the following question:

\begin{quote}
Can a learned, asymmetric feature-interaction
operator support effective regression without
the additional architectural machinery of
a Transformer?
\end{quote}

We study this question through VANTA, a family
of lightweight regression architectures based on
three principles:

\begin{enumerate}[leftmargin=*]
    \item Learn separate affine query and key
    representations for individual input features.

    \item Construct signed feature interactions
    without softmax normalization, using the
    original feature values for aggregation.

    \item Preserve the original feature
    representation through a residual connection,
    followed by a linear regression readout.
\end{enumerate}

We investigate two residual parameterizations,
VANTA-R and VANTA-G, and examine their
mathematical properties and preliminary
empirical performance.

% ============================================================
\section{Related Work}
% ============================================================

\subsection{Feature Interaction Models}

Factorization Machines model pairwise interactions
using low-dimensional latent factors
\cite{rendle2010fm}.

Attentional Factorization Machines extend this
idea by learning interaction-specific weights
through an attention mechanism
\cite{xiao2017afm}.

Deep and Cross Networks provide another approach
to explicitly modeling bounded-degree feature
interactions \cite{wang2017dcn}.

VANTA shares the general motivation of learning
structured interactions between tabular features.

Its specific formulation uses separate affine
query and key projections, direct scalar-value
aggregation, and a residual regression readout.

These connections are important when evaluating
the novelty and contribution of the architecture.

\subsection{Attention and Linear Complexity}

The Transformer introduced scaled dot-product
attention as a general mechanism for relating
input representations \cite{vaswani2017attention}.

Subsequent research demonstrated that certain
attention formulations can exploit the
associativity of matrix multiplication to avoid
explicitly constructing a quadratic attention
matrix \cite{katharopoulos2020linear}.

VANTA uses an unnormalized identity-kernel
dot-product interaction operator, permitting
a similar associative computation.

However, VANTA is not a sequence model and does
not attempt to approximate softmax attention.

\subsection{Tabular Deep Learning}

Gorishniy et al.\ investigated neural-network
architectures for tabular data, including
Transformer-based models and residual networks
\cite{gorishniy2021tabular}.

Grinsztajn et al.\ demonstrated the continued
competitiveness of tree-based models on typical
tabular learning problems
\cite{grinsztajn2022trees}.

These findings motivate comparisons against
both parameter-matched neural baselines and
gradient-boosted decision trees.

% ============================================================
\section{Problem Formulation}
% ============================================================

Consider a supervised regression dataset

\begin{equation}
\mathcal{D}
=
\left\{
(\mathbf{x}^{(n)},y^{(n)})
\right\}_{n=1}^{N},
\end{equation}

where

\begin{equation}
\mathbf{x}^{(n)}\in\mathcal{X},
\qquad
y^{(n)}\in\R.
\end{equation}

Each observation contains $F$ numerical
or categorical features.

Our objective is to learn a function

\begin{equation}
f_\theta:\mathcal{X}\rightarrow\R
\end{equation}

that minimizes the empirical mean squared error

\begin{equation}
\mathcal{L}(\theta)
=
\frac{1}{N}
\sum_{n=1}^{N}
\left(
f_\theta(\mathbf{x}^{(n)})
-y^{(n)}
\right)^2.
\end{equation}

The parameter vector $\theta$ includes
feature embeddings, affine query-key parameters,
residual scaling parameters, and the
final regression readout.

% ============================================================
\section{VANTA Architecture}
% ============================================================

\subsection{Feature Tokenization}

Let $d$ denote the latent interaction dimension.

Each feature is represented by a token

\begin{equation}
\mathbf{t}_i\in\R^d.
\end{equation}

For a standardized numerical feature $x_i$,
we define

\begin{equation}
\mathbf{t}_i=x_i\mathbf{1}_d,
\qquad
v_i=x_i,
\end{equation}

where $\mathbf{1}_d$ is the all-ones vector.

For categorical feature $i$ with category
index $c_i$, we use learned embeddings

\begin{equation}
\mathbf{t}_i=E_i[c_i]\in\R^d
\end{equation}

and a separate learned scalar value embedding

\begin{equation}
v_i=e_i[c_i]\in\R.
\end{equation}

Thus, each original input feature contributes
exactly one token and one scalar value.

Collecting the feature values gives

\begin{equation}
\mathbf{v}
=
\begin{bmatrix}
v_1 & v_2 & \cdots & v_F
\end{bmatrix}^{\top}
\in\R^F.
\end{equation}

This construction preserves feature identity
without requiring one-hot expansion into
additional attention positions.

\subsection{Affine Query and Key Projections}

Each feature receives separate learned
query and key parameters.

Define

\begin{align}
\mathbf{q}_i
&=
\mathbf{t}_i\od\mathbf{w}_i^Q
+
\mathbf{b}_i^Q,
\\
\mathbf{k}_i
&=
\mathbf{t}_i\od\mathbf{w}_i^K
+
\mathbf{b}_i^K,
\end{align}

where

\begin{equation}
\mathbf{w}_i^Q,\mathbf{w}_i^K,
\mathbf{b}_i^Q,\mathbf{b}_i^K
\in\R^d.
\end{equation}

The symbol $\od$ denotes elementwise
multiplication.

For numerical features, these expressions
reduce to conventional affine projections:

\begin{align}
\mathbf{q}_i
&=
x_i\mathbf{w}_i^Q+\mathbf{b}_i^Q,
\\
\mathbf{k}_i
&=
x_i\mathbf{w}_i^K+\mathbf{b}_i^K.
\end{align}

Collecting query and key vectors by rows gives

\begin{equation}
Q,K\in\R^{F\times d}.
\end{equation}

The distinct query and key parameterizations
allow directional interactions between features.

\subsection{Non-Softmax Interaction Operator}

We define the feature-interaction matrix as

\begin{equation}
\boxed{
A(\mathbf{x})
=
\frac{Q(\mathbf{x})K(\mathbf{x})^\top}
{\sqrt d}.
}
\label{eq:attention}
\end{equation}

Its individual elements are

\begin{equation}
A_{ij}
=
\frac{\mathbf{q}_i^\top\mathbf{k}_j}
{\sqrt d}.
\end{equation}

Unlike softmax attention, this formulation
does not normalize rows into probability
distributions.

Consequently, interaction scores may be
positive or negative, and their magnitudes
are not constrained by a simplex normalization.

In general,

\begin{equation}
A_{ij}\neq A_{ji}.
\end{equation}

The interaction matrix is therefore neither
necessarily symmetric nor positive semidefinite.

Moreover, since $A$ is constructed from
a product of two $F\times d$ matrices,

\begin{equation}
\boxed{
\operatorname{rank}(A)\leq d.
}
\label{eq:rank}
\end{equation}

The latent dimension $d$ therefore
constrains the rank of the interaction
operator.

\subsection{Direct Value Aggregation}

Rather than introducing a separate learned
value-projection matrix, VANTA aggregates
the scalar feature values directly:

\begin{equation}
\boxed{
\mathbf{h}
=
A(\mathbf{x})\mathbf{v}.
}
\label{eq:aggregation}
\end{equation}

For feature $i$,

\begin{equation}
h_i
=
\sum_{j=1}^{F}A_{ij}v_j.
\end{equation}

Substituting the interaction scores,

\begin{equation}
h_i
=
\frac{1}{\sqrt d}
\sum_{j=1}^{F}
(\mathbf{q}_i^\top\mathbf{k}_j)v_j.
\end{equation}

Each resulting feature representation
therefore depends on the learned
relationships between the feature's
query and all other feature keys.

The output is

\begin{equation}
\mathbf{h}\in\R^F.
\end{equation}

% ============================================================
\section{Residual Architectures}
% ============================================================

We investigate two methods of combining
the interaction-transformed representation
with the original feature values.

\subsection{VANTA-R: Scalar-Gated Residual Attention}

VANTA-R introduces a learned scalar

\begin{equation}
\gamma\in\R.
\end{equation}

The residual representation is

\begin{equation}
\boxed{
\mathbf{z}_R
=
\mathbf{v}+\gamma\mathbf{h}.
}
\label{eq:residual_r}
\end{equation}

The regression prediction is

\begin{equation}
\boxed{
\hat y_R
=
\mathbf{w}^{\top}\mathbf{z}_R+b.
}
\label{eq:prediction_r}
\end{equation}

Here,

\begin{equation}
\mathbf{w}\in\R^F,
\qquad b\in\R.
\end{equation}

Expanding the expression,

\begin{equation}
\hat y_R
=
\mathbf{w}^{\top}\mathbf{v}
+
\gamma\mathbf{w}^{\top}
A(\mathbf{x})\mathbf{v}
+b.
\end{equation}

The architecture combines a direct
feature pathway with a learned
interaction pathway.

When $\gamma=0$, the model reduces
to a linear readout over the scalar
feature-value representation.

\subsection{VANTA-G: Feature-Wise Gated Residual Attention}

VANTA-G replaces the scalar gate with
a learned vector

\begin{equation}
\boldsymbol{\gamma}\in\R^F.
\end{equation}

The residual representation becomes

\begin{equation}
\boxed{
\mathbf{z}_G
=
\mathbf{v}
+
\boldsymbol{\gamma}\od\mathbf{h}.
}
\label{eq:residual_g}
\end{equation}

The final prediction is

\begin{equation}
\boxed{
\hat y_G
=
\mathbf{w}^{\top}\mathbf{z}_G+b.
}
\label{eq:prediction_g}
\end{equation}

Component-wise,

\begin{equation}
z_{G,i}=v_i+\gamma_i h_i.
\end{equation}

This parameterization permits independently
learned residual scaling for every feature.

No additional feed-forward network or
input-dependent gating network is used.

% ============================================================
\section{Theoretical Properties}
% ============================================================

\subsection{Associative Linear-Time Aggregation}

A naive implementation constructs

\begin{equation}
A=\frac{QK^\top}{\sqrt d}
\end{equation}

and subsequently computes

\begin{equation}
\mathbf{h}=A\mathbf{v}.
\end{equation}

Explicitly constructing the interaction
matrix requires $\mathcal{O}(F^2d)$
arithmetic and $\mathcal{O}(F^2)$ storage.

However, matrix multiplication is associative:

\begin{equation}
(QK^\top)\mathbf{v}
=
Q(K^\top\mathbf{v}).
\end{equation}

Therefore,

\begin{equation}
\boxed{
\mathbf{h}
=
\frac{Q(K^\top\mathbf{v})}{\sqrt d}.
}
\label{eq:linear_attention}
\end{equation}

First compute

\begin{equation}
\mathbf{s}=K^\top\mathbf{v}
\in\R^d.
\end{equation}

Then compute

\begin{equation}
\mathbf{h}
=
\frac{Q\mathbf{s}}{\sqrt d}.
\end{equation}

Both operations require
$\mathcal{O}(Fd)$ arithmetic.

Hence, for scalar feature values,

\begin{equation}
\boxed{
T(F,d)=\mathcal{O}(Fd).
}
\end{equation}

This result eliminates the need to
materialize an $F\times F$ matrix
during ordinary forward inference.

The full matrix can still be constructed
when explicit interaction inspection
is required.

\subsection{Polynomial Interaction Structure}

Consider numerical features exclusively.

Let

\begin{align}
\mathbf{q}_i
&=
x_i\mathbf{a}_i+\mathbf{b}_i,
\\
\mathbf{k}_j
&=
x_j\mathbf{c}_j+\mathbf{e}_j,
\\
v_j&=x_j.
\end{align}

The aggregated representation becomes

\begin{equation}
h_i
=
\frac{1}{\sqrt d}
\sum_j
x_j
(x_i\mathbf{a}_i+\mathbf{b}_i)^\top
(x_j\mathbf{c}_j+\mathbf{e}_j).
\end{equation}

Expanding gives

\begin{align}
h_i=\frac{1}{\sqrt d}\sum_j
\Big[
&
x_ix_j^2\mathbf{a}_i^\top\mathbf{c}_j
\nonumber\\
&+
x_ix_j\mathbf{a}_i^\top\mathbf{e}_j
\nonumber\\
&+
x_j^2\mathbf{b}_i^\top\mathbf{c}_j
\nonumber\\
&+
x_j\mathbf{b}_i^\top\mathbf{e}_j
\Big].
\end{align}

Thus, with affine projections and
identity-kernel attention, the interaction
branch represents a structured polynomial
of degree at most three in numerical inputs.

This includes linear, quadratic, and
cubic contributions.

Categorical embeddings introduce
category-dependent coefficients, while
preserving the bounded polynomial degree
in the numerical inputs.

This limited-degree structure is a
potential inductive bias, but also
restricts the functions representable
by the architecture.

\subsection{Equivalence of Residual Parameterizations}

We now establish a relationship between
VANTA-R and VANTA-G.

\begin{proposition}
Assume that the affine query parameters
are unconstrained.

Then VANTA-R and VANTA-G define the same
family of regression functions for
fixed tokenization and latent dimension.
\end{proposition}

\begin{proof}

For VANTA-G, the $i$-th transformed
feature is

\begin{equation}
z_{G,i}
=
v_i+
\frac{\gamma_i}{\sqrt d}
\mathbf{q}_i^\top
(K^\top\mathbf{v}).
\end{equation}

Define the reparameterized query

\begin{equation}
\tilde{\mathbf{q}}_i
=
\gamma_i\mathbf{q}_i.
\end{equation}

Since

\begin{equation}
\mathbf{q}_i
=
\mathbf{t}_i\od\mathbf{w}_i^Q
+\mathbf{b}_i^Q,
\end{equation}

we have

\begin{equation}
\tilde{\mathbf{q}}_i
=
\mathbf{t}_i\od
(\gamma_i\mathbf{w}_i^Q)
+
\gamma_i\mathbf{b}_i^Q.
\end{equation}

These are valid affine query parameters.

Consequently,

\begin{equation}
z_{G,i}
=
v_i+
\frac{1}{\sqrt d}
\tilde{\mathbf{q}}_i^\top
(K^\top\mathbf{v}),
\end{equation}

which is a VANTA-R representation
with scalar gate equal to one.

Conversely, VANTA-R is a special case
of VANTA-G obtained by setting

\begin{equation}
\gamma_i=\gamma
\qquad\forall i.
\end{equation}

Therefore the two parameterizations
represent the same function family.
\end{proof}

\begin{remark}
Representational equivalence does not
imply identical optimization trajectories.

The two models may exhibit different
gradient dynamics, initialization effects,
and interactions with weight decay.

Their gates are also not uniquely
identifiable as feature-importance measures.
\end{remark}

% ============================================================
\section{Parameter Complexity}
% ============================================================

Let $C$ denote the set of categorical
features, and let $m_i$ be the
embedding-table cardinality of categorical
feature $i$, including the reserved
unknown-category index.

The affine query and key projections require

\begin{equation}
4Fd
\end{equation}

parameters in total.

Categorical token and scalar value
embeddings require

\begin{equation}
(d+1)\sum_{i\in C}m_i
\end{equation}

parameters.

The linear regression readout requires
$F+1$ parameters.

Therefore, the total parameter counts are

\begin{align}
P_R
&=
4Fd
+
(d+1)\sum_{i\in C}m_i
+
F+2,
\\
P_G
&=
4Fd
+
(d+1)\sum_{i\in C}m_i
+
2F+1.
\end{align}

Their difference is

\begin{equation}
\boxed{P_G-P_R=F-1.}
\end{equation}

For the Diamonds configuration,
$F=9$, $d=24$, and
$\sum_{i\in C}m_i=23$.

Hence,

\begin{align}
P_R&=1450,\\
P_G&=1458.
\end{align}

These counts include embedding-table
entries reserved for unknown categories.

% ============================================================
\section{Experimental Setup}
% ============================================================

\subsection{Dataset}

We evaluate the initial architectures
on the Diamonds dataset available
through Kaggle.

The dataset contains 53,940 observations
with nine predictive features and a
continuous price target in USD.

The numerical features are:

\begin{quote}
carat, depth, table, x, y, z.
\end{quote}

The categorical features are:

\begin{quote}
cut, color, clarity.
\end{quote}

Observations with nonpositive prices
or invalid geometric dimensions are
removed using a fixed preprocessing rule.

This leaves 53,917 observations.

\subsection{Preprocessing}

Numerical features are imputed using
training-set medians and standardized
using training-set statistics.

Categorical vocabularies are fitted
using the training split only.

A reserved index represents missing
or previously unseen categories.

The regression target is standardized
using the training-set mean and
standard deviation.

Predictions are transformed back into
USD before metric calculation.

\subsection{Data Splitting}

We use a random split with seed 42:

\begin{itemize}
    \item Training: 43,133 observations.
    \item Validation: 5,392 observations.
    \item Testing: 5,392 observations.
\end{itemize}

Validation error is used for checkpoint
selection and early stopping.

The reported results correspond to
initial single-seed experiments.

\subsection{Training Configuration}

Both VANTA variants are trained using
the following initial configuration:

\begin{center}
\begin{tabular}{ll}
\toprule
Hyperparameter & Value\\
\midrule
Latent dimension & 24\\
Optimizer & AdamW\\
Learning rate & $10^{-3}$\\
Weight decay & $10^{-4}$\\
Batch size & 512\\
Maximum epochs & 200\\
Early-stopping patience & 25\\
Gradient clipping norm & 1.0\\
Initial residual gate & 0.1\\
Loss function & MSE\\
\bottomrule
\end{tabular}
\end{center}

Training was performed using PyTorch
on a CPU.

% ============================================================
\section{Experimental Results}
% ============================================================

\subsection{VANTA-R and VANTA-G}

Table~\ref{tab:variants} reports the
initial results.

\begin{table}[ht]
\centering
\caption{Regression performance of VANTA variants.}
\label{tab:variants}
\begin{tabular}{lrrrr}
\toprule
Model & Parameters & Train $R^2$
& Val. $R^2$ & Test $R^2$\\
\midrule
VANTA-R & 1450 & 0.97415 & 0.97608 & 0.97404\\
VANTA-G & 1458 & 0.97431 & 0.97614 & 0.97421\\
\bottomrule
\end{tabular}
\end{table}

VANTA-G achieves a slightly higher
test $R^2$ than VANTA-R.

However, the difference is small and
cannot establish a reliable advantage
without repeated experiments.

The near-identical results are
consistent with the representational
equivalence established earlier.

\subsection{Comparison with Established Baselines}

We compare VANTA-R against classical
linear regression, polynomial Ridge
regression, a parameter-matched
embedding MLP, XGBoost, and CatBoost.

\begin{table}[ht]
\centering
\caption{Test performance on Diamonds.
Lower RMSE and MAE are better.}
\label{tab:baselines}
\small
\begin{tabular}{lrrr}
\toprule
Model & $R^2$ & RMSE (USD) & MAE (USD)\\
\midrule
MLR & 0.92476 & 1081.91 & 714.82\\
Polynomial Ridge & 0.97126 & 668.66 & 393.11\\
VANTA-R & 0.97404 & 635.49 & 373.60\\
VANTA-G & 0.97421 & 633.44 & 374.02\\
Embedding MLP & 0.97925 & 568.22 & 312.94\\
XGBoost & 0.97928 & 567.76 & 284.23\\
CatBoost & \textbf{0.98227} & \textbf{525.17}
& \textbf{272.27}\\
\bottomrule
\end{tabular}
\end{table}

VANTA-R reduces test RMSE by approximately
41.3\% relative to ordinary multiple
linear regression.

It also outperforms the degree-two
polynomial Ridge baseline on this split.

However, the parameter-matched MLP,
XGBoost, and CatBoost achieve stronger
predictive performance.

Therefore, the results support the
viability of the proposed interaction
operator but do not establish
state-of-the-art regression performance.

\subsection{Attention Ablation Studies}

We conduct ablations of VANTA-R
to investigate the contribution of
learned interactions.

The following variants are evaluated:

\begin{enumerate}[leftmargin=*]
\item Attention disabled and model retrained.
\item Query and key parameters frozen at initialization.
\item Residual gate fixed to $\gamma=1$.
\item Original VANTA-R with learned residual gate.
\end{enumerate}

\begin{table}[ht]
\centering
\caption{VANTA-R ablation results on the test split.}
\label{tab:ablations}
\begin{tabular}{lrr}
\toprule
Configuration & Test $R^2$ & RMSE (USD)\\
\midrule
No attention & 0.92445 & 1084.09\\
Frozen Q/K & 0.92927 & 1048.97\\
Fixed $\gamma=1$ & 0.97394 & 636.76\\
Learned $\gamma$ & \textbf{0.97404} & \textbf{635.49}\\
\bottomrule
\end{tabular}
\end{table}

Disabling attention produces
performance close to the ordinary
linear regression baseline.

Freezing query and key parameters
provides a relatively small improvement
over removing attention entirely.

Training the query and key parameters
produces a substantially larger gain.

These findings support the conclusion
that learned feature interactions are
useful for the observed regression task.

Fixing the scalar gate to one produces
performance nearly identical to
learning it.

This agrees with the theoretical
redundancy of the gate under
unconstrained affine projections.

\subsection{Post-Training Attention Removal}

We additionally disable attention
in the already trained VANTA-R model
without retraining the readout.

The test $R^2$ decreases from
0.97404 to 0.64470.

This indicates substantial
co-adaptation between the trained
interaction pathway and regression
readout.

However, post-training removal
should not be confused with the
performance of an independently
retrained attention-free model.

The latter achieves a test $R^2$
of 0.92445 and is the more
appropriate architectural ablation.

% ============================================================
\section{Discussion}
% ============================================================

The initial results suggest that
a low-rank feature-interaction operator
can provide useful nonlinear structure
for tabular regression.

Three observations are noteworthy.

\paragraph{Learned interactions matter.}

The substantial performance reduction
observed when query-key learning is
disabled suggests that the interaction
branch contributes meaningful predictive
information beyond the direct residual
pathway.

\paragraph{Gate complexity is not essential.}

Scalar and feature-wise residual gating
produce nearly identical accuracy.

This observation is consistent with
their mathematical equivalence under
unconstrained affine query projections.

\paragraph{Nonlinear baselines remain stronger.}

The parameter-matched MLP and
gradient-boosted models outperform VANTA
on the current dataset.

This suggests that the restricted
interaction structure may be less
flexible than the nonlinear functions
learned by these alternatives.

It remains unknown whether VANTA's
structural constraints become advantageous
in lower-data regimes or on datasets
with different feature relationships.

% ============================================================
\section{Limitations}
% ============================================================

The present study has several limitations.

\begin{enumerate}[leftmargin=*]

\item \textbf{Limited dataset coverage.}
The main benchmark currently consists
of one real-world regression dataset.
Cross-dataset evaluation is required.

\item \textbf{Single-seed evidence.}
The reported results do not yet
establish variance across initialization
seeds or data splits.

\item \textbf{Restricted expressivity.}
With identity-kernel attention,
the numerical interaction branch
is polynomial of degree at most three.

\item \textbf{Unnormalized scores.}
Raw dot-product attention permits
signed interactions but can produce
poorly controlled numerical scales.

\item \textbf{Interpretability.}
Attention coefficients and residual
gate magnitudes cannot automatically
be interpreted as causal or unique
feature-importance estimates.

\item \textbf{Benchmark reuse.}
The Diamonds test split has been
examined during iterative model
development, limiting its suitability
as a fully untouched final benchmark.

\item \textbf{Novelty assessment.}
Relationships to factorization machines,
cross networks, and existing
attention-based feature models
require further theoretical comparison.

\end{enumerate}

% ============================================================
\section{Future Work}
% ============================================================

We identify four immediate research directions.

\subsection{Initialization Robustness}

Evaluate VANTA-R, VANTA-G, and
parameter-matched baselines over
multiple random seeds.

Report mean performance, variance,
paired differences, and convergence
characteristics.

\subsection{Cross-Dataset Generalization}

Evaluate on regression datasets
with varying sample sizes,
feature counts, categorical
cardinalities, and interaction
structures.

Of particular interest is whether
the low-rank interaction bias improves
sample efficiency in small-data settings.

\subsection{Normalization and Stability}

Investigate input normalization,
query-key scaling, interaction magnitude,
and alternatives to unrestricted
raw dot-product aggregation.

\subsection{Alternative Interaction Kernels}

Extend the interaction definition to

\begin{equation}
A_{ij}
=
g\left(
\frac{\mathbf{q}_i^\top\mathbf{k}_j}
{\sqrt d}
\right),
\end{equation}

where $g$ may be nonlinear.

Such extensions may increase expressive
capacity, but arbitrary nonlinearities
do not preserve the exact associative
factorization of identity-kernel
attention.

Any resulting accuracy improvements
must therefore be evaluated against
additional computational cost.

% ============================================================
\section{Conclusion}
% ============================================================

We presented VANTA, a lightweight family
of tabular regression architectures
based on learned asymmetric feature
interactions and direct value aggregation.

The architecture removes several
components of conventional Transformers,
including softmax normalization,
learned dense value projections,
feed-forward attention blocks, and
deep stacks of attention layers.

We formalized two variants:
VANTA-R with scalar-gated residual
aggregation and VANTA-G with
feature-wise gated aggregation.

We showed that the identity-kernel
interaction operator admits an
$\mathcal{O}(Fd)$ associative
implementation and that both residual
parameterizations define the same
function family under unconstrained
affine query projections.

Initial experiments demonstrate
strong predictive performance relative
to linear and quadratic regression,
while also identifying stronger
nonlinear baselines.

The present findings motivate broader
experimental and theoretical evaluation
of learned low-rank feature-interaction
operators as an inductive bias
for tabular regression.

% ============================================================
% References
% ============================================================

\begin{thebibliography}{99}

\bibitem{vaswani2017attention}
A. Vaswani et al.
\newblock Attention Is All You Need.
\newblock \emph{NeurIPS}, 2017.

\bibitem{rendle2010fm}
S. Rendle.
\newblock Factorization Machines.
\newblock \emph{IEEE ICDM}, 2010.
\newblock
\url{https://doi.org/10.1109/ICDM.2010.127}

\bibitem{xiao2017afm}
J. Xiao et al.
\newblock Attentional Factorization Machines:
Learning the Weight of Feature Interactions
via Attention Networks.
\newblock \emph{IJCAI}, 2017.
\newblock
\url{https://doi.org/10.24963/ijcai.2017/435}

\bibitem{wang2017dcn}
R. Wang, B. Fu, G. Fu, and M. Wang.
\newblock Deep \& Cross Network for
Ad Click Predictions.
\newblock \emph{ADKDD}, 2017.
\newblock
\url{https://arxiv.org/abs/1708.05123}

\bibitem{katharopoulos2020linear}
A. Katharopoulos, A. Vyas,
N. Pappas, and F. Fleuret.
\newblock Transformers are RNNs:
Fast Autoregressive Transformers
with Linear Attention.
\newblock \emph{ICML}, 2020.
\newblock
\url{https://proceedings.mlr.press/v119/katharopoulos20a.html}

\bibitem{gorishniy2021tabular}
Y. Gorishniy, I. Rubachev,
V. Khrulkov, and A. Babenko.
\newblock Revisiting Deep Learning
Models for Tabular Data.
\newblock \emph{NeurIPS}, 2021.
\newblock
\url{https://arxiv.org/abs/2106.11959}

\bibitem{grinsztajn2022trees}
L. Grinsztajn, E. Oyallon,
and G. Varoquaux.
\newblock Why Do Tree-Based Models
Still Outperform Deep Learning
on Typical Tabular Data?
\newblock \emph{NeurIPS}, 2022.
\newblock
\url{https://arxiv.org/abs/2207.08815}

\end{thebibliography}

\end{document}

